% GENERATED FILE. Edit canonical lessons, then run npm run generate:books.
% canonical-source-hash: fnv1a64:955b482360e91e4e
% canonical-lessons: PA-C305-vadhi, PA-C305-lunna, PA-C305-dati, PA-C305-gahuna, PA-C305-dher

\chapter{Harvest}
\label{ch:pa-a2-harvest}
\hlchaptermodality{305}

\begin{chapteropening}
\textbf{By the end of this chapter:} \emph{I can talk about the harvest.}

Complete the last lesson of chapter 305: I can talk about the harvest.
\end{chapteropening}

\section[\texorpdfstring{\pa{ਵਾਢੀ} (vāḍhī)}{ਵਾਢੀ (vāḍhī)}]{\pa{ਵਾਢੀ} (vāḍhī) — the harvest; reaping}

\label{lesson:PA-C305-vadhi}



\begin{quote}
Before the new one: say the Punjabi for an ear of grain, then the Punjabi for a grain.
\end{quote}

\subsection*{You'll want to know: \pa{ਵਾਢੀ}}
\textbf{\pa{ਵਾਢੀ}} — \emph{vāḍhī} — \textquotedblleft{}the harvest; reaping\textquotedblright{}.

\begin{quote}
\pa{ਵਾਢੀ} \pa{ਵਿਸਾਖ} \pa{ਵਿੱਚ} \pa{ਹੋਵੇਗੀ।}
\end{quote}

\begin{quote}
\emph{vāḍhī visākh vich hovegī.}
\end{quote}

\begin{quote}
The reaping will be in Visakh.
\end{quote}

A feminine noun.

\begin{cousinweb}
\textbf{\pa{ਵਾਢੀ}} (\emph{vāḍhī}) is inherited from Sanskrit \emph{vṛddhi}, \textquotedblleft{}cutting off\textquotedblright{}; CDIAL lists Punjabi \emph{vāḍhī}, \textquotedblleft{}harvesting, harvest\textquotedblright{}, under it (CDIAL 12075).
\end{cousinweb}

\begin{grammarlens}[title={What you've built}]
One more word for this chapter.
\end{grammarlens}

\subsection*{Your turn}
\begin{itemize}
\raggedright
  \item \emph{Say it:} \emph{vāḍhī}
  \item \emph{Say it:} \emph{vāḍhī}, once more
  \item \emph{Say it:} \emph{vāḍhī visākh vich hovegī.} — The reaping will be in Visakh.
  \item \emph{Recall:} say the Punjabi for an ear of grain, then the Punjabi for a grain, then say \emph{vāḍhī} again
\end{itemize}

\begin{tcolorbox}[breakable,colback=teal!4,colframe=teal!35!black,title={Before you move on}]
What does \pa{ਵਾਢੀ} mean? (\textquotedblleft{}The harvest; reaping\textquotedblright{}.) Say it once more.
\end{tcolorbox}
\section[\texorpdfstring{\pa{ਲੁਣਨਾ} (luṇnā)}{ਲੁਣਨਾ (luṇnā)}]{\pa{ਲੁਣਨਾ} (luṇnā) — to reap; to cut a crop}

\label{lesson:PA-C305-lunna}



\begin{quote}
Before the new one: say the Punjabi for a grain, then the Punjabi for reaping.
\end{quote}

\subsection*{You'll want to know: \pa{ਲੁਣਨਾ}}
\textbf{\pa{ਲੁਣਨਾ}} — \emph{luṇnā} — \textquotedblleft{}to reap; to cut a crop\textquotedblright{}.

\begin{quote}
\pa{ਅਸੀਂ} \pa{ਅਗਲੇ} \pa{ਹਫ਼ਤੇ} \pa{ਕਣਕ} \pa{ਲੁਣਾਂਗੇ।}
\end{quote}

\begin{quote}
\emph{asīṁ agle hafte kaṇak luṇāṅge.}
\end{quote}

\begin{quote}
We will reap the wheat next week.
\end{quote}

Past (I …): \pa{ਮੈਂ} \pa{ਕਣਕ} \pa{ਲੁਣੀ}, \emph{maiṁ kaṇak luṇī}.

\begin{cousinweb}
\textbf{\pa{ਲੁਣਨਾ}} (\emph{luṇnā}) is inherited from Sanskrit \emph{lunāti}, \textquotedblleft{}cuts, reaps\textquotedblright{}; CDIAL lists Punjabi \emph{luṇnā}, \textquotedblleft{}cuts, reaps\textquotedblright{}, under it (CDIAL 11082).
\end{cousinweb}

\begin{grammarlens}[title={What you've built}]
One more word for this chapter.
\end{grammarlens}

\subsection*{Your turn}
\begin{itemize}
\raggedright
  \item \emph{Say it:} \emph{luṇnā}
  \item \emph{Say it:} \emph{luṇnā}, once more
  \item \emph{Say it:} \emph{asīṁ agle hafte kaṇak luṇāṅge.} — We will reap the wheat next week.
  \item \emph{Recall:} say the Punjabi for a grain, then the Punjabi for reaping, then say \emph{luṇnā} again
  \item \emph{Write it:} copy \textbf{\pa{ਲੁਣਨਾ}} once, checking the whole word
\end{itemize}

\begin{tcolorbox}[breakable,colback=teal!4,colframe=teal!35!black,title={Before you move on}]
What does \pa{ਲੁਣਨਾ} mean? (\textquotedblleft{}To reap; to cut a crop\textquotedblright{}.) Say it once more.
\end{tcolorbox}
\section[\texorpdfstring{\pa{ਦਾਤੀ} (dātī)}{ਦਾਤੀ (dātī)}]{\pa{ਦਾਤੀ} (dātī) — a sickle}

\label{lesson:PA-C305-dati}



\begin{quote}
Before the new one: say the Punjabi for reaping, then the Punjabi for to reap.
\end{quote}

\subsection*{You'll want to know: \pa{ਦਾਤੀ}}
\textbf{\pa{ਦਾਤੀ}} — \emph{dātī} — \textquotedblleft{}a sickle\textquotedblright{}.

\begin{quote}
\pa{ਕਿਸਾਨ} \pa{ਦਾਤੀ} \pa{ਨਾਲ} \pa{ਕਣਕ} \pa{ਲੁਣਦਾ} \pa{ਹੈ।}
\end{quote}

\begin{quote}
\emph{kisān dātī nāl kaṇak luṇdā hai.}
\end{quote}

\begin{quote}
The farmer reaps the wheat with a sickle.
\end{quote}

A feminine noun.

\begin{cousinweb}
\textbf{\pa{ਦਾਤੀ}} (\emph{dātī}) is inherited from Sanskrit \emph{dātra}, \textquotedblleft{}knife, sickle\textquotedblright{}; CDIAL lists Punjabi \emph{dātī}, \textquotedblleft{}sickle\textquotedblright{}, under it (CDIAL 6260).
\end{cousinweb}

\begin{grammarlens}[title={What you've built}]
One more word for this chapter.
\end{grammarlens}

\subsection*{Your turn}
\begin{itemize}
\raggedright
  \item \emph{Say it:} \emph{dātī}
  \item \emph{Say it:} \emph{dātī}, once more
  \item \emph{Say it:} \emph{kisān dātī nāl kaṇak luṇdā hai.} — The farmer reaps the wheat with a sickle.
  \item \emph{Recall:} say the Punjabi for reaping, then the Punjabi for to reap, then say \emph{dātī} again
  \item \emph{Write it:} copy \textbf{\pa{ਦਾਤੀ}} once, checking the whole word
\end{itemize}

\begin{tcolorbox}[breakable,colback=teal!4,colframe=teal!35!black,title={Before you move on}]
What does \pa{ਦਾਤੀ} mean? (\textquotedblleft{}A sickle\textquotedblright{}.) Say it once more.
\end{tcolorbox}
\section[\texorpdfstring{\pa{ਗਾਹੁਣਾ} (gāhuṇā)}{ਗਾਹੁਣਾ (gāhuṇā)}]{\pa{ਗਾਹੁਣਾ} (gāhuṇā) — to thresh; to tread out grain}

\label{lesson:PA-C305-gahuna}



\begin{quote}
Before the new one: say the Punjabi for to reap, then the Punjabi for a sickle.
\end{quote}

\subsection*{You'll want to know: \pa{ਗਾਹੁਣਾ}}
\textbf{\pa{ਗਾਹੁਣਾ}} — \emph{gāhuṇā} — \textquotedblleft{}to thresh; to tread out grain\textquotedblright{}.

\begin{quote}
\pa{ਬਲਦ} \pa{ਕਣਕ} \pa{ਗਾਹੁੰਦੇ} \pa{ਹਨ।}
\end{quote}

\begin{quote}
\emph{balad kaṇak gāhunde han.}
\end{quote}

\begin{quote}
The oxen thresh the wheat.
\end{quote}

Past (I …): \pa{ਮੈਂ} \pa{ਕਣਕ} \pa{ਗਾਹੀ}, \emph{maiṁ kaṇak gāhī}.

\begin{cousinweb}
\textbf{\pa{ਗਾਹੁਣਾ}} (\emph{gāhuṇā}) is inherited from Sanskrit \emph{gāhatē}, \textquotedblleft{}dives into\textquotedblright{}; CDIAL lists Punjabi \emph{gāhṇā}, \textquotedblleft{}to tread out\textquotedblright{}, under it (CDIAL 4152).
\end{cousinweb}

\begin{grammarlens}[title={What you've built}]
One more word for this chapter.
\end{grammarlens}

\subsection*{Your turn}
\begin{itemize}
\raggedright
  \item \emph{Say it:} \emph{gāhuṇā}
  \item \emph{Say it:} \emph{gāhuṇā}, once more
  \item \emph{Say it:} \emph{balad kaṇak gāhunde han.} — The oxen thresh the wheat.
  \item \emph{Recall:} say the Punjabi for to reap, then the Punjabi for a sickle, then say \emph{gāhuṇā} again
  \item \emph{Write it:} copy \textbf{\pa{ਗਾਹੁਣਾ}} once, checking the whole word
\end{itemize}

\begin{tcolorbox}[breakable,colback=teal!4,colframe=teal!35!black,title={Before you move on}]
What does \pa{ਗਾਹੁਣਾ} mean? (\textquotedblleft{}To thresh; to tread out grain\textquotedblright{}.) Say it once more.
\end{tcolorbox}
\section[\texorpdfstring{\pa{ਢੇਰ} (ḍher)}{ਢੇਰ (ḍher)}]{\pa{ਢੇਰ} (ḍher) — a heap; a pile}

\label{lesson:PA-C305-dher}



\begin{quote}
Before the new one: say the Punjabi for a sickle, then the Punjabi for to thresh.
\end{quote}

\subsection*{You'll want to know: \pa{ਢੇਰ}}
\textbf{\pa{ਢੇਰ}} — \emph{ḍher} — \textquotedblleft{}a heap; a pile\textquotedblright{}.

\begin{quote}
\pa{ਖੇਤ} \pa{ਵਿੱਚ} \pa{ਦਾਣਿਆਂ} \pa{ਦਾ} \pa{ਵੱਡਾ} \pa{ਢੇਰ} \pa{ਹੈ।}
\end{quote}

\begin{quote}
\emph{khet vich dāṇiāṁ dā vaḍḍā ḍher hai.}
\end{quote}

\begin{quote}
There is a big heap of grain in the field.
\end{quote}

A masculine noun.

\begin{cousinweb}
\textbf{\pa{ਢੇਰ}} (\emph{ḍher}) is inherited from a reconstructed \emph{*ḍhēra}, \textquotedblleft{}lump, heap\textquotedblright{}; CDIAL lists Punjabi \emph{ḍher}, \textquotedblleft{}heap\textquotedblright{}, under it (CDIAL 5599).
\end{cousinweb}

\begin{grammarlens}[title={What you've built}]
One more word for this chapter.
\end{grammarlens}

\subsection*{Your turn}
\begin{itemize}
\raggedright
  \item \emph{Say it:} \emph{ḍher}
  \item \emph{Say it:} \emph{ḍher}, once more
  \item \emph{Say it:} \emph{khet vich dāṇiāṁ dā vaḍḍā ḍher hai.} — There is a big heap of grain in the field.
  \item \emph{Recall:} say the Punjabi for a sickle, then the Punjabi for to thresh, then say \emph{ḍher} again
  \item \emph{Write it:} copy \textbf{\pa{ਢੇਰ}} once, checking the whole word
\end{itemize}

\begin{tcolorbox}[breakable,colback=teal!4,colframe=teal!35!black,title={Before you move on}]
What does \pa{ਢੇਰ} mean? (\textquotedblleft{}A heap; a pile\textquotedblright{}.) Say it once more.
\end{tcolorbox}
